% Unicode-support via fontspec (lualatex)
\usepackage{fontspec}
\setmainfont{DejaVu Serif}
\setsansfont{DejaVu Sans}
\setmonofont{DejaVu Sans Mono}

% Verplicht apart voorblad voor reader-PDF's.
% Pandoc's standaard LaTeX-template zet \maketitle direct voor de inhoudsopgave.
% Zonder herdefinitie kunnen titel en TOC op dezelfde pagina terechtkomen.
\makeatletter
\renewcommand{\maketitle}{%
  \begin{titlepage}
    \centering
    \vspace*{0.18\textheight}
    {\Huge\bfseries \@title\par}
    \vspace{1.5em}
    {\Large \@author\par}
    \vfill
    {\large \@date\par}
  \end{titlepage}
}
\makeatother

% Pagebreak na inhoudsopgave
\let\oldtableofcontents\tableofcontents
\renewcommand{\tableofcontents}{\oldtableofcontents\newpage}

% Schaal afbeeldingen automatisch naar maximaal 70% van de paginabreedte.
% Gebruik \includeimage{pad}{breedte} voor afwijkende breedte per afbeelding.
\usepackage{graphicx}
\makeatletter
\def\maxwidth#1{\ifdim\Gin@nat@width>#1#1\else\Gin@nat@width\fi}
\makeatother
\setkeys{Gin}{width=\maxwidth{0.7\linewidth},keepaspectratio}
